\chapter{Zur kirchlichen Chronik}
\quellseite{7}

\section{America}

\subsection{Visitation in Brasilien}
Die letzte Delegatensynode sprach sich dahin aus, daß so bald als möglich durch die Allgemeine Missionscommission eine Visitation in Brasilien vorgenommen werden möchte. Dieser Beschluß der Delegatensynode wird, will’s Gott, dieses Jahr ausgeführt werden. Herr P. L. Lochner von Chicago, ein Glied der Commission, ist mit dem amtlichen Besuch unserer Brüder in Brasilien beauftragt worden. Auch hat die von Herrn P. Lochner bediente Gemeinde aus Liebe zum Reiche Gottes die Erlaubniß gegeben, daß ihr Pastor den ihm gewordenen Auftrag ausrichten kann. So wird denn, will’s Gott, Herr P. Lochner am 5. April von New York nach Brasilien abreisen. F. P.

\subsection{Brasilien}
Das vom Ehrw. Präsidium des Südlichen Districts einer Zeit angeordnete Colloquium mit P. August Zander in Cerro Branco bei Rincao Sao Pedro, der sich zur Aufnahme in unsere Synode gemeldet hatte, ist am 13. November vorigen Jahres durch die Pastoren Mahler, Wittrock und Harder in Gegenwart der Pastoren Kern und Möller abgehalten worden. Die völlige Uebereinstimmung P. Zanders mit unserer Synode in Lehre und Praxis konnte am Schluß des Colloquiums constatirt werden. L. L.

\subsection[Die große Weltausstellung]{Die große Weltausstellung in St. Louis}
Die große Weltausstellung in St. Louis wird am 30. April eröffnet und voraussichtlich auch von vielen Lutheranern besucht werden. Wiederholt sind schon Anfragen namentlich aus der Synodalconferenz, aber auch aus andern lutherischen Kreisen, selbst aus Europa, hier eingelaufen, ob nicht für die die Ausstellung besuchenden Lutheraner passende Quartiere genannt werden könnten. Wir glauben deshalb vielen Lesern des „Lutheraner“ einen Dienst zu thun, wenn wir ihnen schon jetzt mittheilen, daß aus unsern hiesigen Gemeinden zwei Glieder zweckentsprechende und bequemgelegene Herbergen für lutherische Weltausstellungsbesucher einrichten. Die eine befindet sich nur vier Blocks vom Weltausstellungsplatz entfernt, ist ein neues, permanentes Gebäude und bietet Raum für 180 Gäste. Wegen näherer Auskunft wende man sich an den Eigenthümer, Mr. Eckhard Wind, c. o. Mr. H. C. Achenbach, 1704 Market St., St. Louis, Mo. Die andere Herberge befindet sich mehr im Centrum der Stadt, etwa zwanzig Minuten Wegs mit der Straßenbahn vom Weltausstellungsplatz entfernt, und bietet Raum für 60 Gäste. Wegen Auskunft wende man sich an Mr. Wm. Tepe, 2723–2725 Laclede Ave., St. Louis, Mo. Einen halben Block östlich vom hiesigen Centralbahnhof (Union Station) entfernt, in der Apotheke eines andern Lutheraners, des schon genannten Herrn H. C. Achenbach, hat unsere hiesige lutherische Stadtmission ein Informationsbureau eröffnet, in welchem jeder Glaubensbruder sich Auskunft über Lage der Herbergen, über Kirchen, Straßenbahnen und sonstige Ortsverhältnisse der Stadt holen kann. L. F.

\subsection{Die Amerikanische Bibelgesellschaft}
Die Amerikanische Bibelgesellschaft hat im vergangenen Jahre 2,058,298 Bibeln gedruckt oder angekauft, und seit ihrer Gründung im Ganzen 72,000,000 Bibeln ausgegeben. Von den mehr als 2 Millionen Bibeln des letzten Jahres wurden in den Vereinigten Staaten 745,423 verbreitet, und zwar kommen davon allein auf New York 325,735, auf Pennsylvania 135,938, auf Illinois 62,878, auf Wyoming und Arizona nur 56 und 87. Nach den Philippinen wurden 11,774 Exemplare gesandt, aber auch Cuba wurde mit 20,398 nicht vergessen, während Africa 6725, China 1425 und Canada 218 erhielten. — Das sind große Zahlen, sagst du, und dazu kommen noch die ungezählten Tausende, die von andern Gesellschaften und Druckereien ausgegeben werden. Aber was nützt die Bibel im Hause, wenn sie nicht auch gelesen wird? Und daran fehlt es vor allem. L. F.

In der römischen Kirche werden jetzt wieder die üblichen Fastenbriefe erlassen, und derjenige des Erzbischofs von St. Louis liegt vor uns. Zuerst wird darin das Fasten für alle Tage der Fastenzeit mit Ausnahme der Sonntage geboten; nur Eine volle Mahlzeit wird gestattet. Dann werden aber gleich wieder so viele Ausnahmen gemacht, daß es nicht allzu schwer sein wird, die Fasten zu halten. Einmal brauchen diejenigen, die unter 21 Jahren sind, und diejenigen, denen das Fasten gesundheitsschädlich sein könnte, überhaupt nicht zu fasten. Sodann wird der Genuß von Fleisch bei der Hauptmahlzeit an allen Tagen gestattet außer Mittwochs und Freitags, an denen aber Fisch gegessen werden darf. Ferner darf Morgens Kaffee, Thee oder dünne Chocolade genossen werden und Abends ein Essen, das nicht den vierten Theil einer gewöhnlichen Mahlzeit übersteigt. Bei diesem Abendessen ist nur Fleisch verboten, Butter, Eier, Käse, Milch dürfen aus besonderer Vergünstigung des Papstes genossen werden. Auch darf gewechselt und die Hauptmahlzeit Abends gehalten werden. Und endlich gibt es noch eine Vergünstigung des Papstes, die seit 1895 in Kraft ist und zehn Jahre gilt, vermöge welcher Arbeiterfamilien alle Tage außer Freitags und einigen besonderen Tagen Fleisch genießen dürfen. Dafür werden sie ermahnt, aus Dankbarkeit eine andere „Kreuzigung“ des Fleisches zu vollziehen, z.\,B. keine alkoholhaltigen Getränke zu genießen. — Das Reich Gottes besteht nicht in Essen und Trinken, Röm. 14, 17. Wir sollen uns aber auch nach Pauli Wort kein Gewissen machen lassen über Speise oder über Trank, Col. 2, 16. Die Fastengebote, wie die römische Kirche sie alljährlich erläßt, sind Menschengebote, die einer rechten Feier der Passionszeit nur schaden, indem sie die Gewissen irre führen. Die rechte Feier der Passionszeit besteht darin, daß wir uns mit bußfertig gläubigem Herzen in Christi Leiden und Sterben versenken, auf daß wir zur Erkenntniß unserer Sünden und zu dem gewissen Glauben gelangen, daß der Heiland um unserer Missethat willen gelitten hat und gestorben ist, und daß wir nun mit Gott versöhnt sind. Durch solchen wahren Glauben an Christum, aber nicht durch Fasten, Menschenwerk und Ceremonien, haben wir Frieden mit Gott und das ewige Leben. L. F.
